\documentclass[../../../include/open-logic-section]{subfiles}

\begin{document}

\olfileid{sth}{replacement}{finiteaxiomatizability}
\olsection{পরিশিষ্ট: সসীম স্বতঃসিদ্ধায়ন}
প্রতিস্থাপন থেকে কয়েকটি ফল নিয়ে এই
অধ্যায় শেষ করি। প্রথম ফলটি
\citet{Montague1961}-এর। লক্ষ করো,
এটি $\ZF$-এর \emph{ভিতরে} প্রমাণ নয়;
$\ZF$ \emph{সম্পর্কে} প্রমাণ:
\begin{thm}\ollabel{zfnotfinitely}
	$\ZF$ সসীমভাবে স্বতঃসিদ্ধায়িত করা
	যায় না। আরও সাধারণভাবে: $\Th{T}$
	সসীম এবং $\Th{T} \Proves \ZF$ হলে
	$\Th{T}$ অসংগতিপূর্ণ।

	(এখানে আমরা নীরবে কেবল সেই
	প্রথম-ক্রমের বাক্যগুলি বিবেচনা করি,
	যাদের একমাত্র অযৌক্তিক আদিম চিহ্ন
	$\in$। আর $\Th{T} \Proves \ZF$
	দিয়ে বোঝাই: প্রত্যেক $\phi \in \ZF$-এর
	জন্য $\Th{T} \Proves \phi$।)
\end{thm}
\begin{proof}
	এমন একটি সসীম $\Th{T}$ নিই যাতে
	$\Th{T} \Proves \ZF$। তাই $\Th{T}$
	প্রতিফলন প্রমাণ করে, অর্থাৎ
	\olref[sth][replacement][ref]{reflectionschema}।
	$\Th{T}$ সসীম বলে এর স্বতঃসিদ্ধগুলি
	একটি সংযোজন $\theta$ হিসেবে লেখা যায়।
	এই সূত্রে প্রতিফলন প্রয়োগে
	$\Th{T} \Proves \exists \beta(\theta \liff \theta^{V_\beta})$।
	স্পষ্টতই $\Th{T} \Proves \theta$;
	ফলে $\Th{T} \Proves \exists \beta\ \theta^{V_\beta}$।

	এবার $\psi(X)$ দিয়ে সংক্ষেপে বোঝাই:
	\[
		\theta^X \land X\text{ পরিযায়ী} \land (\forall Y \in X)(Y\text{ পরিযায়ী}\lif \lnot \theta^{Y})
	\]
	মোটামুটি এটি বলে: $X$ হলো
	$\theta$-র একটি পরিযায়ী মডেল এবং
	এই ধর্মে $\in$-ক্ষুদ্রতম।
	$\Th{T} \Proves \exists \beta\ \theta^{V_\beta}$
	মনে রেখে, $\ZF$-এ এবং তাই
	$\Th{T}$-তে স্তরাঙ্কের মৌলিক তথ্য
	ব্যবহার করে পাই:
	\begin{equation}
		\Th{T} \Proves \exists M \psi(M). \tag{*}\label{Mpsi}
	\end{equation}
	$\psi(X)$-র প্রথম সংযোজক ব্যবহার
	করে, $\Th{T} \Proves \sigma$ হলে
	$\Th{T} \Proves \forall X(\psi(X) \lif \sigma^X)$।
	তাই \eqref{Mpsi} অনুযায়ী:
	\begin{align*}
		\Th{T} &\Proves \forall X(\psi(X) \lif (\exists N \psi(N))^X)\\
	\intertext{এটি এবং আবার \eqref{Mpsi} ব্যবহার করে:}
		\Th{T} &\Proves \exists M(\psi(M) \land (\exists N \psi(N))^M)
	\intertext{বিশেষত:}
		\Th{T} &\Proves \exists M(\psi(M) \land (\exists N \in M)((N\text{ পরিযায়ী})^N \land (\theta^N)^M))
	\intertext{পরিযায়িতা-সংক্রান্ত প্রাথমিক যুক্তিতে:}
		\Th{T} &\Proves \exists M(\psi(M) \land (\exists N \in M)(N\text{ পরিযায়ী} \land \theta^N))
	\end{align*}
	অতএব $\Th{T}$ অসংগতিপূর্ণ।\footnote{এই
	``প্রাথমিক যুক্তি''তে পরিযায়ী সেটের
	জন্য কয়েকটি ``নিরপেক্ষতা-তথ্য'' প্রমাণ
	করতে হয়।}
\end{proof}

অনুরূপ একটি ফল \citet[223]{Potter2004}
উল্লেখ করেছেন:

\begin{prop}\ollabel{finiteextensionofZ}
	$\Th{T}$-তে $\Z$-এর সঙ্গে সসীম
	সংখ্যক নতুন স্বতঃসিদ্ধ যোগ করা হোক।
	$\Th{T} \Proves \ZF$ হলে
	$\Th{T}$ অসংগতিপূর্ণ। (এখানে
	\olref{zfnotfinitely}-এর মতোই
	নীরব বিধিনিষেধ প্রযোজ্য।)
\end{prop}
\begin{proof}
	পৃথকীকরণ ছকের (অসীম সংখ্যক)
	দৃষ্টান্ত ছাড়া $\Th{T}$-র অন্য
	সব স্বতঃসিদ্ধের সংযোজনকে
	$\theta$ বলি। \olref{zfnotfinitely}-এর
	মতো $\theta$ থেকে $\psi$ সংজ্ঞায়িত
	করে দেখানো যায়,
	$\Th{T} \Proves \exists M \psi(M)$।

	\olref{zfnotfinitely}-এর মতোই
	এমন ছক প্রতিষ্ঠা করা যায় বলে
	উৎসে দাবি করা হয়েছে: যখনই
	$\Th{T} \Proves \sigma$, তখনই
	$\Th{T} \Proves \forall X(\psi(X) \lif \sigma^X)$।
	এরপর ঠিক \olref{zfnotfinitely}-এর
	মতো প্রমাণ শেষ হবে।

	কিন্তু এই ছক প্রতিষ্ঠায়
	\olref{zfnotfinitely}-এর তুলনায়
	আরও কাজ দরকার। কারণ পৃথকীকরণ
	ছকের দৃষ্টান্তগুলি $\Th{T}$-তে
	আছে, কিন্তু $\theta$-র সংযোজক নয়।
	উৎসে বলা হয়েছে, প্রতিটি
	পৃথকীকরণ-দৃষ্টান্ত $\sigma$-র জন্য
	$\Th{T} \Proves \forall X(X\text{ পরিযায়ী} \lif \sigma^X)$
	প্রমাণ করে বাধাটি কাটানো যায়।
	এটি পাঠকের জন্য রাখা হয়েছে।
	% BN-SRC-446: transitivity alone does not imply relativized Separation; the proposed exercise has a finite counterexample.
	\emph{সম্পাদকীয় সতর্কতা: শেষ দাবিটি
	উৎসে যেভাবে লেখা, সেভাবে সত্য নয়।
	শুধু পরিযায়িতা থেকে কোনো সেটের
	সকল উপসেট ওই সেটে থাকবে, তা
	অনুসৃত হয় না; তাই নিচের অনুশীলন
	এবং এই প্রমাণ-ধাপ অসম্পূর্ণ।}
\end{proof}
\begin{prob}
	প্রত্যেক পৃথকীকরণ-দৃষ্টান্ত
	$\sigma$-র জন্য
	$\Z \Proves \forall X(X\text{ পরিযায়ী} \lif \sigma^X)$
	দেখাও। (এই ছক
	\olref[sth][replacement][finiteaxiomatizability]{finiteextensionofZ}-এ
	ব্যবহার করেছি।) সম্পাদকীয় সতর্কতাটি
	মনে রেখো: এই রূপে দাবিটির
	একটি প্রতিউদাহরণ আছে।
\end{prob}
\begin{prob}
	প্রত্যেক $\phi \in \Z$-এর জন্য
	$\ZF \Proves \phi^{V_{\omega+\omega}}$
	দেখাও।
\end{prob}
\begin{prob}
	\olref[sth][replacement][finiteaxiomatizability]{zfnotfinitely}
	ও \olref[sth][replacement][finiteaxiomatizability]{finiteextensionofZ}-এর
	প্রমাণে ব্যবহৃত বাকি ছকগুলি
	যাচাই করো।
\end{prob}

\olref[sth][replacement][absinf]{sec}-এ
যেমন বলেছি, এটি দেখায় যে প্রতিস্থাপন
\olref[ordinals][ordtype]{thmOrdinalRepresentation}-এর
চেয়ে কঠোরভাবে শক্তিশালী। আরও
সুনির্দিষ্টভাবে: $\Z$ + ``প্রত্যেক
সুক্রমবিন্যাস একটি অনন্য ক্রমসংখ্যার
সমরূপ'' সঙ্গতিপূর্ণ হলে,
এ থেকে প্রতিস্থাপনের কোনো একটি
দৃষ্টান্ত প্রমাণ করা যায় না।


	%	By assumption, $\psi^{V_\alpha}$ has the form:
	%	\[
	%		(\forall A \in V_\alpha)(\exists S \in V_\alpha)(\forall x \in V_\alpha)(x \in S \liff (\phi^{V_\alpha}(x) \land x \in A))
	%	\]
	%	To establish this holds, fix $A \in V_\alpha$. Using Separation, obtain:
	%	\[
	%		S = \Setabs{x \in A}{\phi^{V_\alpha}(x)}
	%	\]
	%	Now $S \in V_\alpha$, since $S \subseteq A \in V_\alpha$, and clearly $(\forall x \in V_\alpha)(x \in S \liff (\phi^{V_\alpha}(x) \land x \in A)$.

\end{document}
